\section{Conclusion}

We have presented a framework for understanding the architectural evolution of autonomous agentic systems, grounded in empirical observations from the Wallfacer experiment and extended through formal and theoretical analysis. The central empirical finding, bottleneck migration, connects to established results in operations research, optimization theory, and cybernetics. The fission principle provides a candidate explanation for observed patterns of role differentiation. The formalization of trust as Gaussian Process preference learning offers a principled approach to the HITL boundary problem. The theoretical extensions to populations, multi-principal scenarios, and self-referential self-models map a large and largely unexplored conceptual space.

We conclude with what we consider the framework's most important, if most speculative, implication. If the pattern of bottleneck migration is genuine and persistent, then the long-term trajectory of agentic architecture is not toward a fixed optimal design but toward a sustained capacity for redesign. The questions practitioners and researchers pose about these systems (``How should I structure the agents?'' at one level; ``How should the system restructure itself?'' at another; ``What constitutes the system's boundary?'' at a third) themselves recapitulate the detect-fission-emerge cycle that the framework describes. Whether this self-similarity is a deep structural feature of complex adaptive systems or merely a suggestive analogy is, we believe, one of the more interesting open questions in the field.
